% Options for packages loaded elsewhere
\PassOptionsToPackage{unicode}{hyperref}
\PassOptionsToPackage{hyphens}{url}
%
\documentclass[
  11pt,
]{article}
\usepackage{amsmath,amssymb}
\usepackage{iftex}
\ifPDFTeX
  \usepackage[T1]{fontenc}
  \usepackage[utf8]{inputenc}
  \usepackage{textcomp} % provide euro and other symbols
\else % if luatex or xetex
  \usepackage{unicode-math} % this also loads fontspec
  \defaultfontfeatures{Scale=MatchLowercase}
  \defaultfontfeatures[\rmfamily]{Ligatures=TeX,Scale=1}
\fi
\usepackage{lmodern}
\ifPDFTeX\else
  % xetex/luatex font selection
\fi
% Use upquote if available, for straight quotes in verbatim environments
\IfFileExists{upquote.sty}{\usepackage{upquote}}{}
\IfFileExists{microtype.sty}{% use microtype if available
  \usepackage[]{microtype}
  \UseMicrotypeSet[protrusion]{basicmath} % disable protrusion for tt fonts
}{}
\makeatletter
\@ifundefined{KOMAClassName}{% if non-KOMA class
  \IfFileExists{parskip.sty}{%
    \usepackage{parskip}
  }{% else
    \setlength{\parindent}{0pt}
    \setlength{\parskip}{6pt plus 2pt minus 1pt}}
}{% if KOMA class
  \KOMAoptions{parskip=half}}
\makeatother
\usepackage{xcolor}
\usepackage[margin=1in]{geometry}
\usepackage{color}
\usepackage{fancyvrb}
\newcommand{\VerbBar}{|}
\newcommand{\VERB}{\Verb[commandchars=\\\{\}]}
\DefineVerbatimEnvironment{Highlighting}{Verbatim}{commandchars=\\\{\}}
% Add ',fontsize=\small' for more characters per line
\newenvironment{Shaded}{}{}
\newcommand{\AlertTok}[1]{\textcolor[rgb]{1.00,0.00,0.00}{\textbf{#1}}}
\newcommand{\AnnotationTok}[1]{\textcolor[rgb]{0.38,0.63,0.69}{\textbf{\textit{#1}}}}
\newcommand{\AttributeTok}[1]{\textcolor[rgb]{0.49,0.56,0.16}{#1}}
\newcommand{\BaseNTok}[1]{\textcolor[rgb]{0.25,0.63,0.44}{#1}}
\newcommand{\BuiltInTok}[1]{\textcolor[rgb]{0.00,0.50,0.00}{#1}}
\newcommand{\CharTok}[1]{\textcolor[rgb]{0.25,0.44,0.63}{#1}}
\newcommand{\CommentTok}[1]{\textcolor[rgb]{0.38,0.63,0.69}{\textit{#1}}}
\newcommand{\CommentVarTok}[1]{\textcolor[rgb]{0.38,0.63,0.69}{\textbf{\textit{#1}}}}
\newcommand{\ConstantTok}[1]{\textcolor[rgb]{0.53,0.00,0.00}{#1}}
\newcommand{\ControlFlowTok}[1]{\textcolor[rgb]{0.00,0.44,0.13}{\textbf{#1}}}
\newcommand{\DataTypeTok}[1]{\textcolor[rgb]{0.56,0.13,0.00}{#1}}
\newcommand{\DecValTok}[1]{\textcolor[rgb]{0.25,0.63,0.44}{#1}}
\newcommand{\DocumentationTok}[1]{\textcolor[rgb]{0.73,0.13,0.13}{\textit{#1}}}
\newcommand{\ErrorTok}[1]{\textcolor[rgb]{1.00,0.00,0.00}{\textbf{#1}}}
\newcommand{\ExtensionTok}[1]{#1}
\newcommand{\FloatTok}[1]{\textcolor[rgb]{0.25,0.63,0.44}{#1}}
\newcommand{\FunctionTok}[1]{\textcolor[rgb]{0.02,0.16,0.49}{#1}}
\newcommand{\ImportTok}[1]{\textcolor[rgb]{0.00,0.50,0.00}{\textbf{#1}}}
\newcommand{\InformationTok}[1]{\textcolor[rgb]{0.38,0.63,0.69}{\textbf{\textit{#1}}}}
\newcommand{\KeywordTok}[1]{\textcolor[rgb]{0.00,0.44,0.13}{\textbf{#1}}}
\newcommand{\NormalTok}[1]{#1}
\newcommand{\OperatorTok}[1]{\textcolor[rgb]{0.40,0.40,0.40}{#1}}
\newcommand{\OtherTok}[1]{\textcolor[rgb]{0.00,0.44,0.13}{#1}}
\newcommand{\PreprocessorTok}[1]{\textcolor[rgb]{0.74,0.48,0.00}{#1}}
\newcommand{\RegionMarkerTok}[1]{#1}
\newcommand{\SpecialCharTok}[1]{\textcolor[rgb]{0.25,0.44,0.63}{#1}}
\newcommand{\SpecialStringTok}[1]{\textcolor[rgb]{0.73,0.40,0.53}{#1}}
\newcommand{\StringTok}[1]{\textcolor[rgb]{0.25,0.44,0.63}{#1}}
\newcommand{\VariableTok}[1]{\textcolor[rgb]{0.10,0.09,0.49}{#1}}
\newcommand{\VerbatimStringTok}[1]{\textcolor[rgb]{0.25,0.44,0.63}{#1}}
\newcommand{\WarningTok}[1]{\textcolor[rgb]{0.38,0.63,0.69}{\textbf{\textit{#1}}}}
\setlength{\emergencystretch}{3em} % prevent overfull lines
\providecommand{\tightlist}{%
  \setlength{\itemsep}{0pt}\setlength{\parskip}{0pt}}
\setcounter{secnumdepth}{5}
\ifLuaTeX
  \usepackage{selnolig}  % disable illegal ligatures
\fi
\usepackage{bookmark}
\IfFileExists{xurl.sty}{\usepackage{xurl}}{} % add URL line breaks if available
\urlstyle{same}
\hypersetup{
  pdftitle={Ordinary zero-Hessian recognition on strict analytic curves over arbitrary fields},
  pdfauthor={P06.2 integration research draft},
  hidelinks,
  pdfcreator={LaTeX via pandoc}}

\title{Ordinary zero-Hessian recognition on strict analytic curves over
arbitrary fields}
\author{P06.2 integration research draft}
\date{4 October 2026}

\usepackage{xurl}
\begin{document}
\maketitle

\subsection{Result and conventions}\label{result-and-conventions}

This document gives a complete integration argument in the curve scope
using the supplied exact supporting theorems n700 and n701 and the
stated primary foundations. The new integration is an unverified draft,
with an author audit. The independent receipts for the supporting
theorems are preserved at their own scopes; no earlier verdict certifies
this integration. Original P06.2 remains open in arbitrary dimension, so
its status is \textbf{partial} and \texttt{candidate=false}. P05.2 is
not addressed.

All spaces are actual Berkovich analytic spaces, with their specified
structures. A curve means a separated analytic space pure of analytic
dimension one, as in Gubler--Rabinoff (GR) 3.1. An ordinary open is open
in the underlying topology. Strictness always means strict analyticity
over the current ground field. Put \(D=d'd''\). All tests and currents
in the initial equation are \textbf{ordinary} Chambert--Loir--Ducros
(CLD) tests and currents. On a curve, \[
 D[h](\psi)=\int_X h(D\psi)_{\mathrm{Rad},K},
 \qquad \psi\in A_c^{0,0}(X).                         \tag{1}
\] Canonical integration retains all actual generic Artinian lengths of
\(X\), including unequal lengths on different components and nilpotent
structure. The smooth top form \(D\psi\) has a signed Radon measure. We
do not assume that \(D[h]\) itself is a Radon measure.

GR-subharmonicity means GR Definition 3.4; GR-harmonicity means both
signs are subharmonic (Definition 3.13). Classical pluriharmonicity
means both signs are classically psh in Definition 4.1 (Definition 7.1).
These definitions include extension fields and curve maps. Real PL below
has the \textbf{G-local} convention of GR 2.16 and GJR Definition 5.1,
on the G-topology generated by strict analytic domains. GJR-harmonicity
means Definition 7.2, equivalently the numerical triviality criterion in
Remark 7.3(1). It is a different regularity assertion from ordinary CLD
smooth harmonicity.

\textbf{Theorem (arbitrary-field intrinsic curve recognition).} Let
\(K\) be any complete nontrivially valued non-Archimedean field. Let
\(X\) be a good, separated, strict, boundaryless, pure one-dimensional
\(K\)-analytic space. Every finite continuous \(h:X\to\mathbb R\)
satisfying (1) equal to zero for every ordinary compact smooth scalar
test is GR-harmonic and classically pluriharmonic. It is real PL in the
stated G-local convention, and GJR-harmonic. The result holds on every
ordinary open subspace of \(X\).

More explicitly, for \textbf{every} analytic extension field \(L/K\),
\textbf{every} \(L\)-analytic curve \(Y\), and \textbf{every}
\(L\)-analytic morphism \(f:Y\to X_L\), the finite continuous function
\[
                         h\circ\pi_{L/K}\circ f           \tag{2}
\] is GR-harmonic. No smoothness, strictness, boundarylessness,
reducedness, or separability hypothesis is added to these test curves or
maps beyond being an analytic curve and morphism in the published
definition. If \(Y\) is smooth and strict, (2) is also
Thuillier-harmonic. The ground-field theorem includes imperfect and
positive-characteristic fields, singular curves, arbitrary nilpotents,
and unequal generic lengths. There is no global algebraic-model or
properness hypothesis.

\subsection{Exact support and immutable
bindings}\label{exact-support-and-immutable-bindings}

We use only Theorem A of n700 for recognition and Theorems A and B of
n701 for completed-algebraic transport and its application checks. Their
complete frozen proof bodies were supplied in the input and are copied
unchanged into this workspace. Historical note labels inside those
bodies refer to the original campaigns. They are not new native support
IDs.

The n700 controlling full argument is the published round07 body
\path{/workspace/research/deliverables/round-07/P06.2-curve-recognition-argument.md},
SHA-256
\texttt{8c46bccc1c2dd4286d3ccfc42aea70698ef8e449b9e57f962d707ecffd953a04}.
Its adjacent published review binds this expanded argument and records
the independent exact-scope PASS with 38 premises and 66 passages in
round06/P06.2-curve-recognition-embedded-evidence-review. It separately
binds the historical main argument with hash
\texttt{d83f011d8165f248c3927ef26591dd791f91573e309e3921396e543afa5bb380}.
We preserve the frozen self-audit wording and the later review record;
the self-audit wording does not override that later scoped receipt.

The n701 controlling \textbf{main} is
\path{/workspace/research/campaigns/round-21/P06.2-all-dimensional-completed-algebraic-continuation/artifacts/codex-ZwIWCf/P06.2-argument.md},
SHA-256
\texttt{e7c350a209fa4b40875fa15a65875d210c306919a089cdc7b8cf918c895dbd42}.
The adjacent published round22 review receipt explicitly binds this main
and its declared proof file, and records the independent exact-scope
PASS with 23 premises and 32 passages in
round21/P06.2-all-dimensional-completed-algebraic-review. The
\textbf{expanded review wrapper}, also preserved, has the different
SHA-256
\texttt{f1577176fe748c3ee6ddd18e3feae02221aba4d4ed60d7acaadb6f8dde93d7e6}.
These files are not identified with one another. The old draft's
unverified-status wording remains immutable; the scoped published
receipt governs its supporting acceptance, without extending its
theorem.

Hashes establish file identity. They do not establish mathematics. The
supporting results used here are the following exact mathematical gates.

\textbf{Gate R (n700, Theorem A).} Over an algebraically closed complete
nontrivially valued field, a finite continuous ordinary-zero-Hessian
function on any good separated strict boundaryless pure analytic curve
is GR-harmonic, classically pluriharmonic, G-locally real PL, and
GJR-harmonic. This includes singular and nonreduced curves, and holds on
every ordinary open. It retains canonical integration with actual
generic lengths. It grants no arbitrary-field current ascent.

\textbf{Gate T (n701, Theorems A and B).} If \(K\) is complete and
nontrivially valued, \(M/K\) is an algebraic valued extension, and
\(C=\widehat M\), then every good separated strict boundaryless pure
\(n\)-dimensional space, \(n\geq1\), has actual scalar extension \(X_C\)
with those same properties and pure dimension \(n\). For finite
continuous \(h\), \[
                   D[h]=0\quad\Longleftrightarrow\quad
                   D[h\circ\pi_{C/K}]=0.                 \tag{3}
\] The equations use all ordinary compact \((n-1,n-1)\) tests and each
field's canonical integration. Every compact ordinary form in any
bidegree on \(X_C\) is the exact smooth pullback of one compact ordinary
form on \(X_E\) for a finite embedded \(K\subseteq E\subseteq M\). These
assertions hold on every ordinary open and retain all actual structures
and lengths. No perfectness, separability, tameness, reducedness, or
preimposed upstairs equation is required. The theorem does not assert
(3) for an arbitrary analytic extension outside this completed-algebraic
scope.

In n700 the smooth recognition step uses actual ordinary interiors of
strict affinoids embedded into smooth proper algebraic curves over the
same algebraically closed field. Wanner comparison is applied only
there. The supplied closed-curve, reduction, puncture-removal, and
normalization proofs then handle singular and nonreduced structures. We
use the intrinsic Theorem A, not its separate ground-field ambient-probe
theorem. In particular, no assertion about normalization smoothness over
the original imperfect \(K\) is imported.

In n701, actual finite tensors, finite Galois induction, Frobenius for
mixed inseparable stages, and exact rational domains supply the finite
and completed gates. Its componentwise purity proof is separate from
mere invariance of total dimension. The proof's ordinary compact gluing
and continuous-multiplier bounds are needed; a scalar-invariant test
claim or a direct-image formula alone would not supply (3).

\subsection{The valued completed algebraic
closure}\label{the-valued-completed-algebraic-closure}

\textbf{Standard finite-field foundation.} A finite algebraic extension
\(E/K\) of a complete non-Archimedean field admits a unique extending
absolute value, is complete in that absolute value, and every embedding
of such finite extensions over \(K\) is isometric. This includes
inseparable extensions. This is BGR, \emph{Non-Archimedean Analysis},
3.2.4/2, used at exactly this scope in n700, S2, and proved in n701,
Section 3.1. The uniqueness assertion follows by finite-dimensional norm
equivalence and applying a comparison bound to arbitrarily high powers.
We use this standard foundation on finite fields over the complete
fields \(K\) and, in the next lemma, \(L\); it is not a current theorem.

\textbf{Lemma 1.} An algebraic closure \(\overline K\) has a compatible
absolute value extending that of \(K\), and \(C=\widehat{\overline K}\)
is an algebraically closed complete nontrivially valued analytic
extension of \(K\), in every characteristic.

\textbf{Proof.} Give each finite intermediate \(E/K\) its unique norm.
Two choices agree on intersections because their restrictions extend the
norm on a finite complete intermediate field. On their finite compositum
they agree with its unique norm. These norms therefore define a
multiplicative ultrametric absolute value on the union \(\overline K\).
The embedding of \(K\) is isometric. Its metric completion is a field:
for a nonzero limit \(a\) of elements \(a_m\) of \(\overline K\),
eventually \(|a_m|=|a|>0\), and \[
 |a_m^{-1}-a_j^{-1}|=|a_m-a_j|/(|a_m||a_j|)
\] shows that the inverses converge. Multiplication and the norm extend
continuously, retaining multiplicativity and the ultrametric inequality.
Completeness follows by construction. An element of \(K^\times\) with
absolute value different from one retains that value in \(C\), so the
valuation is nontrivial. This verifies the analytic-extension definition
in GR 2.1--2.2: a complete non-Archimedean field with an isometric
embedding of the ground field.

Here is a proof of algebraic closure, rather than an invocation of a
completion theorem. Take a monic polynomial \[
 P(T)=T^d+\sum_{i<d}a_iT^i\in C[T],\qquad d\geq1,
 \quad R=\max(1,\max_{i<d}|a_i|).
\] For \(m\geq1\), choose monic \(P_m\in\overline K[T]\) whose
coefficients have errors at most \(\delta_m=2^{-md}R^{-(d-1)}\) from
those of \(P\). Density allows these simultaneous choices. Every root of
\(P_m\) has norm at most \(R\): if \(|z|>R\), the leading term has norm
strictly greater than every other term. All \(P_m\) split, with
repetitions allowed, in \(\overline K\).

Choose a root \(b_1\) of \(P_1\). Given a root \(b_m\) of \(P_m\), \[
 |P_{m+1}(b_m)|\leq\delta_m R^{d-1}=2^{-md}.
\] Factor \(P_{m+1}(T)=\prod_{j=1}^d(T-c_j)\). At least one factor has
\(|b_m-c_j|\leq2^{-m}\); otherwise their product would have norm greater
than \(2^{-md}\). Choose it as \(b_{m+1}\). The ultrametric inequality
gives \(|b_j-b_m|\leq2^{-m}\) for \(j>m\), so \(b_m\) converges in \(C\)
to \(b\). Coefficient convergence and the bound \(|b_m|\leq R\) give
\(P(b)=\lim_m P_m(b_m)=0\). Thus every nonconstant polynomial has a
root, after dividing by its leading coefficient; repeated division gives
algebraic closure. No derivative, separability, discrete valuation, or
spherical completeness was used. \(\square\)

\textbf{Lemma 2 (actual common field for this integration).} For every
analytic extension \(L/K\) there is an actual complete nontrivially
valued field \(F\) with isometric embeddings of \(C\) and \(L\) whose
restrictions to \(K\) agree.

\textbf{Proof.} Choose an algebraic closure \(\overline L\) and endow it
with the compatible extension norm using the complete field \(L\) as in
Lemma 1. A \(K\)-embedding \(\overline K\to\overline L\) exists:
partially defined embeddings on intermediate fields, ordered by
extension, have unions on chains. A maximal one missing an algebraic
element extends by choosing a root of that element's minimal polynomial
in \(\overline L\), a contradiction. Each finite \(K\)-subfield embeds
isometrically by the finite-field foundation. Thus this embedding is
isometric on all of \(\overline K\). Take \(F=\widehat{\overline L}\).
The embedding extends isometrically from \(C\) into this complete field,
and \(L\) already embeds isometrically. They agree on \(K\).
Nontriviality is inherited from \(K\). \(\square\)

Lemma 2 is a direct construction for the fields appearing here. It does
not assert that \(C\widehat\otimes_K L\) is a field or define any
pullback of a continuous current. It can be used in the function-level
scalar descent below. In that use we make no claim that \(F/C\) is
completed-algebraic, and no current equation is transported through it.

\subsection{Published function-level inputs and their
applications}\label{published-function-level-inputs-and-their-applications}

All GR references below refer to \textbf{arXiv:2506.13548v2, 17
September 2025}, W. Gubler and J. Rabinoff. The inspected primary text
has SHA-256
\texttt{3a221b86116604c950db07e4c60ab184701f0624deebabd92502e5c81871436b};
the PDF has SHA-256
\texttt{83ca9a19394f44bb2e4e1aa316de95938f311acd53a00fe250fa0f605184b726}.
The operator-selected bodies are extracted by LF alone, preserving
formfeeds: GR 3.5(7) at LF 827--885; 3.14--3.15 and 3.16 at 1062--1103;
4.1--4.2 at 1105--1177. The first unit starts with ``Proof of
Proposition 3.5(7)'' and ends with its printed proof square; the last
includes the complete printed proof of 4.2, not just its statement.
Their exact body hashes are, respectively,
\texttt{c2133ed53c731ad28c6d3049f47e3f3af27023f374e57fe0c06dac3570866455},
\texttt{6faf0cbe62ec3b5cfb3d2817b457a27ad9922b80938b455d944eace90e79f793},
and
\texttt{3aaee5628a226898b74714bb8b2e15108bf61d7d0e8d3b96afeea80f7377446f}.
The primary appendix in the authoritative argument embeds these full
units and the additional definitions used below. No source ending is
replaced by a tool-output truncation.

\textbf{GR subharmonicity and harmonicity.} Definition 3.4 starts with
an upper semicontinuous function on a separated pure analytic curve
\(Z\). For every nontrivially valued analytic extension \(E/k\) and
every map from a smooth connected strict \(E\)-curve into \(Z_E\), its
composite is Thuillier-subharmonic or identically \(-\infty\).
Definition 3.13 says a real-valued function is harmonic if both signs
are subharmonic. GR 3.4.2 and 3.12 identify these finite two-sign
functions on smooth strict curves over nontrivially valued fields with
Thuillier harmonic functions. A finite function never uses the
identically-\(-\infty\) alternative.

\textbf{GR Proposition 3.14(6).} Precisely: if \(Z\) is a \(k\)-analytic
curve, \(u:Z\to\mathbb R\) is continuous, and \(E/k\) is an analytic
extension field, then \(u\) is harmonic if and only if
\(u\circ\pi_{E/k}:Z_E\to\mathbb R\) is harmonic. There is no assumption
of reducedness, smoothness, strictness, perfectness, characteristic, or
algebraicity of \(E/k\). This is the two-sign specialization of
Proposition 3.5(7), for upper semicontinuous functions. We apply it
first with \(Z=X\) and \(E=C\), \textbf{after} Gate R recognizes
\(h_C\).

We have checked the complete scalar-extension proof, rather than reading
``if and only if'' as current functoriality. Its reverse direction uses
a simultaneous analytic extension and, on a smooth strict probe, the
comparison definition of Thuillier subharmonicity. In our application
the simultaneous extension of \(C\) and any other \(L/K\) is explicitly
constructed in Lemma 2. The proof then compares already subharmonic
functions, never continuous currents.

The exact smooth foundations in that proof are: Thuillier Definition
3.1.5 compares an upper semicontinuous function, not identically
\(-\infty\) on any component, with harmonic functions on every strict
affinoid domain; Thuillier Proposition 2.3.18 preserves the harmonic
sheaf on a smooth strict curve under any non-Archimedean extension; and
GR Lemma 2.12 states for \textbf{any} analytic space \(U\) and analytic
extension \(E/k\) that \[
                         \pi_{E/k}(\partial U_E)=\partial U.       \tag{4}
\] Scalar projection is surjective (GR 2.9, CLD 11.1.9). The affinoid
harmonic pullback used in GR's comparison step has the same continuous
boundary values, via the strict-affinoid harmonic convention and
Thuillier Lemma 2.3.5; the printed proof applies Proposition 2.3.18 at
precisely this step. These assertions concern harmonic functions and
relative boundaries, not an equality of all analytic interiors.

For clarity, the comparison step is as follows. On the smooth strict
probe \(Y\) over \(L\), take a strict affinoid \(U\subseteq Y\) and a
harmonic \(a\) on \(U\) with \(v\leq a\) on \(\partial U\). After
passing to the common field, the recognized subharmonic function \(v_F\)
and harmonic \(a_F\) satisfy the same boundary inequality by (4).
Thuillier comparison gives \(v_F\leq a_F\) on \(U_F\). Surjectivity then
gives \(v\leq a\) on \(U\). Each sign gives the required two-sign
property. Smoothness and strictness are those of the probe in Definition
3.4; their preservation under complete scalar extension is the standard
base-change property of quasi-smoothness (CLD 11.4.8--9), together with
boundarylessness. If componentwise purity needs to be checked for the
particular \(Y_F\) used here, \(F=\widehat{\overline L}\) is exactly an
instance of Gate T over \(L\), applied \textbf{only to geometry}, not to
any current. Finite values exclude the \(-\infty\) component alternative
in the printed proof. Thus the source's use of component surjectivity to
exclude that alternative is immaterial for these applications.

This reconstructs the applicable function-level logic of the published
scalar detection result. We import that result in its exact published
scope; we do not extrapolate the arbitrary-current gate to \(F/C\).

\textbf{GR curve pullback and classical psh.} Proposition 3.14(7) states
that for any morphism of analytic curves \(f:Y\to Z\), harmonic \(u\)
pulls back to harmonic \(u\circ f\); its converse requires \(f\) finite
surjective. We use only forward curve pullback in the new integration.
Proposition 3.5(8) gives the corresponding forward subharmonic pullback.
Definition 4.1 requires, for every analytic extension \(L/k\), every
\(L\)-analytic curve \(Y\), and every \(f:Y\to Z_L\), subharmonicity of
\(u\circ\pi_{L/k}\circ f\). GR 4.1.2 proves that on an analytic curve
classical psh is equivalent to subharmonicity, using precisely 3.5(7,8).
It covers all analytic test curves, not just the smooth strict probes
used to define subharmonicity. We apply it separately to \(h\) and
\(-h\).

GR Proposition 4.2(5), also fully supplied, is a second valid route:
classical psh functions on \textbf{any} analytic space are preserved and
detected by every analytic field extension. Once Gate R has made both
signs of \(h_C\) classically psh, it too descends the same conclusion.
It is a consistency check here, not an additional current premise. We do
not use the finite-map converse in 4.2(6), although its full source
proof is retained in the selected unit.

\textbf{GR Proposition 3.15 and GJR.} Precisely: on a strict analytic
curve over a nontrivially valued field, a continuous function is
GR-harmonic if and only if it is real PL and its residual line bundle
\(L_h(x)\in\operatorname{Pic}(\operatorname{red}(X,x))_{\mathbb R}\) is
numerically trivial for every \(x\). The statement has no reducedness or
perfectness requirement. Its full printed proof uses completed
perfection, reduction, and normalization to reduce to the rig-smooth
case, and descends already recognized harmonicity and PL data. We apply
this \textbf{published proposition itself}, not a claim that its
intermediate currents ascend. Our \(K\) is nontrivially valued, \(X\) is
strict and good, and \(h\) is continuous and now harmonic, so all its
hypotheses hold. GR 2.16 fixes the G-local analytic-unit convention.

The GJR primary version is \emph{Forms on Berkovich spaces based on
harmonic tropicalizations}, arXiv:2111.05741v3, 7 March 2025, the
inspected final-published-version snapshot. Its field convention 2.2
requires complete nontrivial valuation. Section 5 and Definition 5.1
concern good strict analytic spaces and G-covers by strict affinoid
domains, with finite sums \(\sum\lambda_j\log|f_j|\) of analytic units.
Definition 7.2 defines harmonicity of a real-PL function by both-sign
semipositivity at every point. Remark 7.3(1) equates that condition with
numerical triviality of \(L_h(x)\). Our ground-field space satisfies
these exact field, goodness, and strictness hypotheses. Proposition 3.15
first supplies the real-PL condition, and only then do we apply the GJR
criterion. No ordinary smooth moment representation follows from this.

\subsection{Actual geometry, compact tests, and calibration at the
current
gate}\label{actual-geometry-compact-tests-and-calibration-at-the-current-gate}

Choose \(C\) as in Lemma 1, with \(M=\overline K\). This is an actual
algebraic valued extension and its metric completion, so Gate T applies
with \(n=1\). On a strict affinoid \(V=\mathcal M(A)\) the source
upstairs is \(\mathcal M(A\widehat\otimes_K C)\), with the actual
completed algebra, including nilpotents.

Goodness follows from inverse images of ordinary interiors of affinoid
neighbourhoods; strict presentations remain strict after scalar
extension. Separatedness follows by base change of the closed diagonal.
Boundarylessness follows from the relative-interior base-change
inclusion, CLD 11.4.1 or Berkovich 1993, Proposition 1.5.5(iii),(iv), as
used in Gate T. We use only that inclusion, not equality of arbitrary
absolute interiors. Purity of \textbf{every} actual upstairs component
and local dimension one follow from n701's separate componentwise purity
proof in Sections 3.2, 6 and 9. Merely knowing \(\dim X_C=\dim X=1\)
would be insufficient.

The function \(h_C=h\circ\pi_{C/K}\) is finite and continuous. A compact
set downstairs has compact inverse image: take a finite cover by
ordinary interiors of affinoid neighbourhoods, whose scalar inverse
images are compact; the inverse image of the compact set is closed in
their finite union since the carriers are Hausdorff. This is the
compact-support assertion in n701 Section 3.3, with no finite-fibre
assumption on \(C/K\).

For an arbitrary \(\beta\in A_c^{0,0}(X_C)\), n701 Theorem A gives a
finite embedded \(E\subseteq\overline K\) and
\(\alpha_E\in A_c^{0,0}(X_E)\) with \(\beta=q_E^*\alpha_E\) \textbf{as
smooth sections}, on the entire space. This includes zero regions,
ordinary presentation opens and compact zero extensions. Its smooth
Hessian is therefore exactly \(q_E^*D\alpha_E\). Finite-stage transport
in n701 gives \(D[h_E]=0\). Canonical smooth-top-measure projection, CLD
Proposition 15.3.1, gives \[
 \begin{aligned}
 \int_{X_C}h_C(D\beta)_{\mathrm{Rad},C}
 &=\int_{X_C}(h_E\circ q_E)(q_E^*D\alpha_E)_{\mathrm{Rad},C}\\
 &=\int_{X_E}h_E(D\alpha_E)_{\mathrm{Rad},E}=0.          \tag{5}
 \end{aligned}
\] The multiplier is bounded on the compact support. This is an
application of a theorem proved for all compact tests. It is not a claim
that an arbitrary continuous current has been pulled back.

The length bookkeeping is essential. In the calibration underlying (5),
a finite Artinian moment-fibre algebra \(B\) over a complete field \(H\)
has dimension \(d\). After a complete extension \(H'/H\), its tensor is
already complete as a finite-dimensional vector space and decomposes
into local factors \(B_j\). A composition series in each factor gives \[
 d=\dim_{H'}(B\otimes_H H')
   =\sum_j\operatorname{length}_{B_j}(B_j)[\kappa_j:H'].     \tag{6}
\] Each composition factor is the residue field of its local Artinian
ring, so this calculation permits nilpotents, inseparable residue
extensions, and unequal lengths. It is the actual degree sum in CLD
13.3.3 and 14.3.6, used in n701 Section 10, and supplies the canonical
projection formula with \textbf{no field-degree factor}. Individual
generic weights may change under scalar extension; their appropriate
calibrated sums are retained. They are never normalized to one or
divided out.

In contrast, CLD Remark 14.4.3 concerns a finite \(E\)-analytic space
viewed as \(K\)-analytic on the \textbf{same carrier}: its
\(K\)-calibration is \([E:K]\) times the \(E\)-calibration. n701 uses
that different statement only after constructing an actual local
\(E\)-analytic structure in its finite-Galois induction. We do not
replace (5) by that statement, and never assign an infinite degree
factor to \(C/K\).

\subsection{Proof of the theorem and all
quantifiers}\label{proof-of-the-theorem-and-all-quantifiers}

If \(X\) is empty the conclusions are vacuous. Otherwise take the valued
completed algebraic closure \(C\) of Lemma 1. The original \(K\) and
\(X\) satisfy every hypothesis of Gate T: complete nontrivial valuation,
goodness, separatedness, strictness, boundarylessness, and pure
dimension one. The preceding application checks give a good separated
strict boundaryless \textbf{pure analytic curve} \(X_C\), with its
actual structure, and the equation \(D[h_C]=0\) against every ordinary
compact scalar test.

The field \(C\) is algebraically closed, complete, and nontrivially
valued. Gate R now applies directly to \(X_C\) and its finite continuous
multiplier. It gives GR-harmonicity of \(h_C\). Its singular and
nonreduced cases are part of that exact theorem, so no smoothness or
reduction assumption on \(X_C\) is needed and none is silently
introduced. In particular, possible new nilpotents from geometric
inseparability cause no scope failure.

Apply GR Proposition 3.14(6) with \(Z=X\), \(k=K\), \(k'=C\), and
\(u=h\). \(X\) is a separated pure analytic curve, \(h\) is continuous,
and \(C/K\) is an analytic extension. The proposition detects
harmonicity of this \textbf{already recognized} upstairs function and
gives GR-harmonicity of \(h\). This is precisely the published reverse
function-level implication; it does not detect or ascend an arbitrary
current.

Consequently \(h\) and \(-h\) are subharmonic in Definition 3.4. GR
4.1.2 applies on the analytic curve \(X\) to each sign, making both
signs classically psh. Definition 7.1 gives classical pluriharmonicity.
To spell out its full scope, fix any analytic extension \(L/K\), any
\(L\)-analytic curve \(Y\), and any \(f:Y\to X_L\). Definition 4.1 for
each sign gives subharmonicity on \(Y\) of the corresponding sign of
(2). The composite is finite and continuous; Definition 3.13 makes it
GR-harmonic. This includes the further extension-field and smooth-probe
quantifiers already inside subharmonicity on \(Y\). If \(Y\) is smooth
and strict, GR 3.4.2 and 3.12 identify it with Thuillier harmonicity. No
current equation on \(X_L\), on \(Y\), or on any common extension has
been used in this quantifier step. The field \(L\) may be transcendental
over \(K\), and the map \(f\) may be inseparable, nonfinite, or
constant.

Now apply GR Proposition 3.15 on the original strict \(K\)-curve. Its
nontrivial-valuation hypothesis and continuity hypothesis hold, and
GR-harmonicity was just established. It gives real PL in the published
G-local convention and numerical triviality of every residual class. GJR
Definition 7.2 and Remark 7.3(1), with the complete nontrivially valued
\(K\) and good strict \(X\), give GJR-harmonicity.

Finally let \(\Omega\subseteq X\) be an ordinary open. It is an analytic
domain (CLD 11.1.7). It is good, separated and strict by ordinary-local
restriction and the strict-neighbourhood statement CLD 11.7.7; local
dimension is preserved by domain restriction (11.2.3), so it is pure
one-dimensional if nonempty. Its open inclusion has full relative
interior, and composition with the boundaryless structural morphism
makes it boundaryless (11.4.1). A compact scalar test on \(\Omega\)
extends by zero to a compact ordinary scalar test on \(X\): its support
is a compact closed subset of the Hausdorff \(X\) contained in the open
\(\Omega\), and every exterior point has a neighbourhood disjoint from
that support. Canonical integration and differentials restrict to
domains (CLD 14.3.5, 15.7.2 and Section 7.7), so (1) restricts. Apply
the same theorem to \(\Omega\), or use harmonicity's sheaf property GR
3.14(1). These are actual ordinary opens at every analytic point, not
merely one affinoid or an admissible cover. This proves all claims.
\(\square\)

\subsection{Comparison with the original residual
scope}\label{comparison-with-the-original-residual-scope}

The catalogue's published curve recognition case was an ordinary open in
a smooth proper algebraic curve's analytification over an algebraically
closed complete nontrivially valued field. Gate R already extends that
to all good separated strict boundaryless pure analytic curves over that
field. The new integration removes the algebraic-closure restriction
\textbf{for every such actual analytic curve}, including imperfect-field
and singular or nonreduced cases. It needs no product, invariance,
factorization, imposed scalar-stability, finite covering presentation,
or algebraic compactification downstairs.

This is a new stated special scope of original P06.2, not its full
arbitrary-dimensional completion. No original-scope analytic
counterexample is constructed. The all-dimensional Gate T supplies
current ascent only through a completion of an algebraic extension; Gate
R supplies recognition only in dimension one. Their composition does not
recognize arbitrary higher-dimensional nonproduct germs. Ordinary opens
of a curve are covered; arbitrary nonproduct higher-dimensional domains
remain outside this conclusion. The non-strict invisible-boundary
example does not meet this theorem's strict hypotheses.

The two original proof obligations remain distinct. The current-ascent
obligation used here is exactly the completed-algebraic gate. The
recognition obligation is then closed by intrinsic curve recognition and
published harmonic-function descent. An ordinary CLD-smooth decreasing
psh approximation of either sign is a stronger construction, and no such
construction is supplied or needed by this route. We make no general
positive-current regularization claim and no ordinary smooth
harmonic-approximation claim. No new arbitrary-transcendental current
ascent, reverse-current characterization, or all-extension zero-current
theorem is claimed. Harmonicity of every field-extended curve composite
is a function-level conclusion.

No round23 incomplete/nullproof, old inconclusive arbitrary-field
recognition, or old full-curve two-sign approximation assertion is a
dependency. The independent round25 smoothing campaign is neither read
as support nor modified. The quoted 4.2(6) finite-map converse is not
used to bypass any missing image or component hypothesis.

\subsection{Evidence, author audit, and
reproduction}\label{evidence-author-audit-and-reproduction}

The authoritative file is \texttt{P06.2-argument.md}, consisting of this
entire core plus the exact selected primary appendix. All new proofs and
application checks are in it. The complete supporting bodies will be
embedded as n700 and n701 by the queued review helper; local copies
\texttt{sources/n700-argument.md} and \texttt{sources/n701-argument.md}
preserve them without edits. Their receipts, freezes, review runs and
bindings are also copied. \texttt{sources/manifest.json} records their
read-only origins and exact hashes. None of the historical status
metadata is rewritten.

\texttt{author-audit.md} checks every gate and conclusion.
\texttt{research-notes.md} records exclusions and remaining gaps.
\texttt{status.json} has exactly the five requested keys, with id P06.2,
status partial, and a nonnull authoritative proof file.
\texttt{notes.json} contains canonical new native notes with ID-only
support arrays. \texttt{P06.2-partial.tex} is a standalone paper with
the complete mathematical core, and \texttt{P06.2-partial.pdf} is built
in two installed pdfLaTeX passes. The primary-source appendix retains
extracted glyph defects; it is evidence, not a replacement for the
precise mathematical statements above.

Rerun from this fresh workspace:

\begin{Shaded}
\begin{Highlighting}[]
\FunctionTok{bash}\NormalTok{ reproduce.sh}
\end{Highlighting}
\end{Shaded}

The rerun copies only read-only sources into this workspace, verifies
bindings, re-extracts the primary units by LF, assembles the argument,
checks schemas and source hashes, generates the standalone TeX with
installed pandoc, compiles it twice, and records outputs and environment
versions. Build and integrity checks establish reproducibility and
typesetting, not mathematical correctness. No extra model, subagent,
network, install, external message, authentication or environment
change, old-artifact write, database operation, or heavy computation is
used. No supervised heavy-compute facility was supplied or needed.

\subsection{References}\label{references}

\begin{enumerate}
\def\labelenumi{\arabic{enumi}.}
\tightlist
\item
  W. Gubler and J. Rabinoff, \emph{Classically psh and pluriharmonic
  functions on Berkovich spaces}, arXiv:2506.13548v2, 17 September 2025.
  Field conventions 2.1--2.2, scalar spaces 2.9, boundaries 2.11--2.12,
  real PL 2.16, curves 3.1, Definitions 3.4 and 3.13, Proposition
  3.5(7,8), 3.4.2, 3.12, Propositions 3.14(1,6,7) and 3.15, Definition
  4.1, 4.1.2, Proposition 4.2(5), Definition 7.1.
\item
  W. Gubler, P. Jell and J. Rabinoff, \emph{Forms on Berkovich spaces
  based on harmonic tropicalizations}, arXiv:2111.05741v3, 7 March 2025,
  inspected final published snapshot. Field convention 2.2, Definition
  5.1, Definition 7.2 and Remark 7.3(1).
\item
  A. Chambert-Loir and A. Ducros, \emph{Formes differentielles reelles
  et courants sur les espaces de Berkovich}, arXiv:1204.6277v2, 25 July
  2025. Sections 7.7, 11.1.7--9, 11.2.3, 11.4.1, 11.4.8--9, 11.7.7,
  13.3.3, 14.3.5--6, 14.4.3, 15.3.1 and 15.7.2.
\item
  A. Thuillier, \emph{Theorie du potentiel sur les courbes en geometrie
  analytique non archimedienne}, thesis, 2005, HAL tel-00010990v1. Lemma
  2.3.5, Proposition 2.3.18, Definition 3.1.5, in the exact
  function-level use in GR's scalar-subharmonicity proof.
\item
  S. Bosch, U. Guentzer and R. Remmert, \emph{Non-Archimedean Analysis},
  Springer, 1984, 3.2.4/2, the finite complete-field norm theorem.
\item
  V. Berkovich, \emph{Etale cohomology for non-Archimedean analytic
  spaces}, Publ. Math. IHES 78 (1993), Proposition 1.5.5(iii),(iv), as
  stated and applied in the full supplied n701 and CLD 11.4.1.
\item
  Supplied n700, exact intrinsic curve recognition, complete round07
  proof and round06 exact-scope review, at the bindings given above.
\item
  Supplied n701, exact completed-algebraic compact-test/current
  transport and componentwise purity, authoritative round21 main and
  round22 exact-scope review, at the distinct bindings given above.
\end{enumerate}

\end{document}
